\documentclass[aspectratio=169,xcolor=dvipsnames]{beamer}
\usetheme{SimpleDarkBlue}
\usepackage{hyperref}
\usepackage{graphicx} % Allows including images
\usepackage{booktabs} % Allows the use of \toprule, \midrule and \bottomrule in tables
\usepackage[spanish]{babel}
\usepackage[utf8]{inputenc}
\usepackage{array}
\usepackage{tabularx}
\usepackage{caption}
\usepackage{adjustbox}
\usepackage{amssymb}
\usepackage{tcolorbox}
\usepackage{amsmath}

\graphicspath{{./graficas}} 

\title{Estadística aplicada con R}
\subtitle{Rmarkdown}

\author{Práxedis Jiménez Ruvalcaba}

\institute
{

}
\date{\today}

\tcbuselibrary{theorems}

% Custom colored theorem box
\newtcbtheorem{mytheo}{Teorema}%
{colback=orange!5!white,colframe=orange!75!black,fonttitle=\bfseries}{th}

\begin{document}

\begin{frame}[c]
    \begin{beamercolorbox}[sep=8pt,center,rounded=true,shadow=true]{title}
        \usebeamerfont{title}\inserttitle\par
        \ifx\insertsubtitle\@empty
        \else
          \vskip0.25em
          {\usebeamerfont{subtitle}\insertsubtitle\par}
        \fi 
    \end{beamercolorbox}

    \vspace{1cm}
    \begin{center}
        \usebeamerfont{author}\insertauthor\par
        \vspace{0.5cm}
        \usebeamerfont{institute}\insertinstitute\par
        \usebeamerfont{date}\insertdate
    \end{center}
    
\end{frame}

\begin{frame}{¿Qué es rmarkdown?}
    \begin{block}{Lenguaje de marcado}
        Sistema de anotación que usa etiquetas para organizar y dar estructura a un documentos digitales, como pudieran
        ser datos (XML), páginas web (HTML), textos legibles con formato sencillo (Markdown), textos y presentaciones 
        con formatos complejos (Latex), etc.
    \end{block}
\end{frame}

\begin{frame}{Notas de uso de Rmarkdown}
    \begin{itemize}
        \item Para generar un archivo rmarkdown, hay que crear un archivo .Rmd
        \item R studio tiene integrado una guía de rmarkdown, a través de la sección \textit{help}.
        \item Cuenta con soporte nativo par exportación de \textit{html}, pero con paquetes puede exportarse como \textit{pdf} o \textit{word}
        \item Recordar parámetro \textit{r, echo = FALSE}, para no mostrar código. Si es lo que se busca.
        \item La sintáxis es prácticamente a Markdown usual, excepto por el código dinámico.
        \item Archivos markdown no toman las variables globales. 
    \end{itemize}
\end{frame}

\begin{frame}[fragile]{Código dinámico}
    Para añdir código y resultados en un bloque, se usa:
    \begin{verbatim}
        ```{r, echo = TRUE}
            código de R
        ```
    \end{verbatim}
    Para añadir resultados de código de manera narrativa, se puede escribir
    \begin{verbatim}
        Tuvimos `r nrow(data)` casos en cuenta.
    \end{verbatim}
    Sin r, se interpreta como un bloque de código genérico (formato verbatim).
\end{frame}

\begin{frame}{Símbolos más usados}
    \begin{tabular}{|c|p{8cm}|}
            \hline
            \textbf{Símbolo} & \textbf{Uso} \\ \hline
            \# & \# Título principal, \#\# Título, \#\#\# Subtítulo \\ \hline
            * * & *itálica*, **negrita** \\ \hline
            * & * ítem de lista \\ \hline
            1. & 1. ítem numerado \\ \hline
            - & más de 3 ----, representan un salto de página\\ \hline
            \$ & \$ecuación\$, \$\$ecuación en bloque\$\$ \\ \hline       
    \end{tabular}
\end{frame}

\begin{frame}{Añadido de tablas}
    Para representar tabla se usa lo siguiente. Se añaden columnas manteniendo consistencia.
    \begin{verbatim}
        
        Primer encabezado  | Segundo encabezado

        ------------- | -------------

        contenido de celda  | contenido de celda

        contenido de celda  | contenido de celda
    \end{verbatim}
\end{frame}

\begin{frame}[fragile]{Referencias e imágenes}
    Para añadir un link a una frase y que quedé como referencia, se usa normalmente
    \begin{verbatim}
        [frase][id]. 

        --al final del documento--
        [id]: http://example.com/ "Título"
    \end{verbatim}
    Para añadir una imagen con una referencia, se utiliza.
    \begin{verbatim}
        ![alt text][id]

        --al final del documento--
        [id]: figures/img.png "Title"
    \end{verbatim}
\end{frame}
\end{document}